%%%%%%%%%%%%%%%%%%%%%%%%%%%%%%%%%%%%%%%%%%%%%%%%%%%%%%%%%%%%
% 10.16 OPTIMAL CONTROL
% The Composition of Advanced Mathematics
%%%%%%%%%%%%%%%%%%%%%%%%%%%%%%%%%%%%%%%%%%%%%%%%%%%%%%%%%%%%

\section{Optimal Control}
\label{sec:optimal-control}

\prerequisite{
Optimization fundamentals, multivariable calculus, differential equations,
dynamical systems, linear algebra, calculus of variations, dynamic
programming, and basic probability for stochastic extensions.
}

\learningobjectives{
By the end of this section, the reader should be able to:
\begin{itemize}
    \item define an optimal-control problem;
    \item distinguish states, controls, trajectories, and parameters;
    \item formulate continuous-time dynamical constraints;
    \item distinguish running and terminal costs;
    \item recognize fixed and free terminal conditions;
    \item understand admissible controls;
    \item distinguish open-loop and feedback control;
    \item connect optimal control with calculus of variations;
    \item formulate the Hamiltonian;
    \item state Pontryagin's Maximum Principle under a consistent sign
          convention;
    \item derive costate equations;
    \item interpret transversality conditions;
    \item understand pointwise Hamiltonian minimization;
    \item identify bang-bang controls;
    \item recognize switching functions;
    \item understand singular controls conceptually;
    \item formulate linear-quadratic optimal-control problems;
    \item derive the finite-horizon Riccati differential equation;
    \item understand linear-quadratic regulation;
    \item recognize infinite-horizon algebraic Riccati equations;
    \item formulate discrete-time optimal-control problems;
    \item derive Bellman equations;
    \item understand value functions;
    \item state the Hamilton--Jacobi--Bellman equation;
    \item connect the HJB equation with feedback control;
    \item distinguish PMP and dynamic-programming approaches;
    \item formulate state and control constraints;
    \item understand endpoint constraints;
    \item recognize minimum-time control;
    \item recognize minimum-energy control;
    \item understand controllability as a feasibility issue conceptually;
    \item understand terminal penalties and tracking objectives;
    \item formulate trajectory-tracking problems;
    \item recognize model predictive control;
    \item distinguish direct and indirect numerical methods;
    \item understand control parameterization and direct transcription;
    \item recognize shooting methods;
    \item understand collocation conceptually;
    \item recognize stochastic optimal-control problems;
    \item understand stochastic HJB equations conceptually;
    \item recognize connections with reinforcement learning;
    \item connect optimal control with robotics, aerospace, engineering,
          economics, energy systems, biological systems, and finance.
\end{itemize}
}

%%%%%%%%%%%%%%%%%%%%%%%%%%%%%%%%%%%%%%%%%%%%%%%%%%%%%%%%%%%%
\subsection{What Is Optimal Control?}
\label{subsec:what-is-optimal-control}
%%%%%%%%%%%%%%%%%%%%%%%%%%%%%%%%%%%%%%%%%%%%%%%%%%%%%%%%%%%%

Optimal control studies how a dynamical system should be influenced over
time in order to optimize a performance criterion.

A typical system is described by

\[
\boxed{
\dot{x}(t)
=
f
(
x(t),
u(t),
t
),
}
\]

where

\[
x(t)
=
\text{state},
\]

and

\[
u(t)
=
\text{control}.
\]

The control is selected to optimize an objective while the state must obey
the differential equation.

%%%%%%%%%%%%%%%%%%%%%%%%%%%%%%%%%%%%%%%%%%%%%%%%%%%%%%%%%%%%
\subsection{State Variables}
\label{subsec:optimal-control-state}
%%%%%%%%%%%%%%%%%%%%%%%%%%%%%%%%%%%%%%%%%%%%%%%%%%%%%%%%%%%%

The state vector

\[
\boxed{
x(t)
\in
\mathbb{R}^n
}
\]

contains the information needed to describe the system at time \(t\).

Examples include:

\[
\boxed{
\text{position and velocity},
}
\]

\[
\boxed{
\text{inventory levels},
}
\]

\[
\boxed{
\text{population sizes},
}
\]

\[
\boxed{
\text{temperature},
}
\]

or

\[
\boxed{
\text{capital stock}.
}
\]

%%%%%%%%%%%%%%%%%%%%%%%%%%%%%%%%%%%%%%%%%%%%%%%%%%%%%%%%%%%%
\subsection{Control Variables}
\label{subsec:optimal-control-control}
%%%%%%%%%%%%%%%%%%%%%%%%%%%%%%%%%%%%%%%%%%%%%%%%%%%%%%%%%%%%

The control

\[
\boxed{
u(t)
\in
\mathbb{R}^m
}
\]

represents quantities the decision maker can manipulate.

Examples include:

\[
\boxed{
\text{thrust},
}
\]

\[
\boxed{
\text{steering input},
}
\]

\[
\boxed{
\text{production rate},
}
\]

\[
\boxed{
\text{drug dosage},
}
\]

or

\[
\boxed{
\text{investment rate}.
}
\]

%%%%%%%%%%%%%%%%%%%%%%%%%%%%%%%%%%%%%%%%%%%%%%%%%%%%%%%%%%%%
\subsection{State Dynamics}
\label{subsec:optimal-control-dynamics}
%%%%%%%%%%%%%%%%%%%%%%%%%%%%%%%%%%%%%%%%%%%%%%%%%%%%%%%%%%%%

The state evolves according to

\[
\boxed{
\dot{x}
=
f(x,u,t).
}
\]

Given an initial condition

\[
\boxed{
x(t_0)
=
x_0,
}
\]

and an admissible control \(u(t)\), the differential equation determines a
state trajectory.

Thus control decisions cannot choose states independently at each time.

%%%%%%%%%%%%%%%%%%%%%%%%%%%%%%%%%%%%%%%%%%%%%%%%%%%%%%%%%%%%
\subsection{Admissible Controls}
\label{subsec:admissible-controls}
%%%%%%%%%%%%%%%%%%%%%%%%%%%%%%%%%%%%%%%%%%%%%%%%%%%%%%%%%%%%

A control is admissible if it satisfies all required conditions.

For example,

\[
\boxed{
u(t)\in U
}
\]

for almost every \(t\), where \(U\) is the control constraint set.

A simple bound is

\[
\boxed{
u_{\min}
\leq
u(t)
\leq
u_{\max}.
}
\]

Admissibility may also require measurability, integrability, or additional
regularity.

%%%%%%%%%%%%%%%%%%%%%%%%%%%%%%%%%%%%%%%%%%%%%%%%%%%%%%%%%%%%
\subsection{The Bolza Problem}
\label{subsec:bolza-problem}
%%%%%%%%%%%%%%%%%%%%%%%%%%%%%%%%%%%%%%%%%%%%%%%%%%%%%%%%%%%%

A standard finite-horizon optimal-control problem is

\[
\boxed{
\min_{u(\cdot)}
J[u]
=
\Phi
(
x(T),
T
)
+
\int_{t_0}^{T}
L
(
x(t),
u(t),
t
)
\,dt
}
\]

subject to

\[
\boxed{
\dot{x}(t)
=
f
(
x(t),
u(t),
t
),
}
\]

and

\[
\boxed{
x(t_0)=x_0.
}
\]

This is called the Bolza form.

%%%%%%%%%%%%%%%%%%%%%%%%%%%%%%%%%%%%%%%%%%%%%%%%%%%%%%%%%%%%
\subsection{Running Cost}
\label{subsec:running-cost}
%%%%%%%%%%%%%%%%%%%%%%%%%%%%%%%%%%%%%%%%%%%%%%%%%%%%%%%%%%%%

The function

\[
\boxed{
L(x,u,t)
}
\]

is the running cost or stage cost.

It accumulates cost continuously over time.

Examples include:

\[
\boxed{
\text{fuel consumption},
}
\]

\[
\boxed{
\text{tracking error},
}
\]

\[
\boxed{
\text{energy use},
}
\]

or

\[
\boxed{
\text{operating cost}.
}
\]

%%%%%%%%%%%%%%%%%%%%%%%%%%%%%%%%%%%%%%%%%%%%%%%%%%%%%%%%%%%%
\subsection{Terminal Cost}
\label{subsec:terminal-cost}
%%%%%%%%%%%%%%%%%%%%%%%%%%%%%%%%%%%%%%%%%%%%%%%%%%%%%%%%%%%%

The function

\[
\boxed{
\Phi(x(T),T)
}
\]

is the terminal cost.

It may penalize:

\[
\boxed{
\text{final position error},
}
\]

\[
\boxed{
\text{remaining inventory},
}
\]

\[
\boxed{
\text{terminal wealth},
}
\]

or other final-state outcomes.

%%%%%%%%%%%%%%%%%%%%%%%%%%%%%%%%%%%%%%%%%%%%%%%%%%%%%%%%%%%%
\subsection{Lagrange and Mayer Forms}
\label{subsec:lagrange-mayer-forms}
%%%%%%%%%%%%%%%%%%%%%%%%%%%%%%%%%%%%%%%%%%%%%%%%%%%%%%%%%%%%

If

\[
\Phi=0,
\]

the problem is in Lagrange form:

\[
\boxed{
J
=
\int_{t_0}^{T}
L(x,u,t)\,dt.
}
\]

If

\[
L=0,
\]

the problem is in Mayer form:

\[
\boxed{
J
=
\Phi(x(T),T).
}
\]

Bolza form contains both.

%%%%%%%%%%%%%%%%%%%%%%%%%%%%%%%%%%%%%%%%%%%%%%%%%%%%%%%%%%%%
\subsection{Fixed and Free Terminal State}
\label{subsec:fixed-free-terminal-state}
%%%%%%%%%%%%%%%%%%%%%%%%%%%%%%%%%%%%%%%%%%%%%%%%%%%%%%%%%%%%

A terminal state may be fixed:

\[
\boxed{
x(T)=x_T.
}
\]

Alternatively, \(x(T)\) may be free and penalized through

\[
\Phi(x(T),T).
\]

Different terminal conditions lead to different transversality conditions.

%%%%%%%%%%%%%%%%%%%%%%%%%%%%%%%%%%%%%%%%%%%%%%%%%%%%%%%%%%%%
\subsection{Fixed and Free Terminal Time}
\label{subsec:fixed-free-terminal-time}
%%%%%%%%%%%%%%%%%%%%%%%%%%%%%%%%%%%%%%%%%%%%%%%%%%%%%%%%%%%%

The final time \(T\) may also be fixed or optimized.

If \(T\) is free, then time itself becomes part of the optimization
problem.

Minimum-time control is an important example.

%%%%%%%%%%%%%%%%%%%%%%%%%%%%%%%%%%%%%%%%%%%%%%%%%%%%%%%%%%%%
\subsection{Open-Loop Control}
\label{subsec:open-loop-control}
%%%%%%%%%%%%%%%%%%%%%%%%%%%%%%%%%%%%%%%%%%%%%%%%%%%%%%%%%%%%

An open-loop control is specified as a function of time:

\[
\boxed{
u=u(t).
}
\]

Once the control trajectory is determined, it does not explicitly respond
to deviations in the observed state.

%%%%%%%%%%%%%%%%%%%%%%%%%%%%%%%%%%%%%%%%%%%%%%%%%%%%%%%%%%%%
\subsection{Feedback Control}
\label{subsec:feedback-control}
%%%%%%%%%%%%%%%%%%%%%%%%%%%%%%%%%%%%%%%%%%%%%%%%%%%%%%%%%%%%

A feedback control uses the current state:

\[
\boxed{
u(t)
=
\pi
(
x(t),
t
).
}
\]

The function

\[
\pi
\]

is called a feedback policy.

Feedback can react to disturbances or modeling errors because the control
depends on the observed state.

%%%%%%%%%%%%%%%%%%%%%%%%%%%%%%%%%%%%%%%%%%%%%%%%%%%%%%%%%%%%
\subsection{Open Loop Versus Feedback}
\label{subsec:open-loop-versus-feedback}
%%%%%%%%%%%%%%%%%%%%%%%%%%%%%%%%%%%%%%%%%%%%%%%%%%%%%%%%%%%%

Open-loop control specifies:

\[
\boxed{
\text{what to do as a function of time}.
}
\]

Feedback control specifies:

\[
\boxed{
\text{what to do as a function of the current state}.
}
\]

Dynamic programming naturally produces feedback policies.

Pontryagin-based boundary-value methods often produce open-loop optimal
trajectories.

%%%%%%%%%%%%%%%%%%%%%%%%%%%%%%%%%%%%%%%%%%%%%%%%%%%%%%%%%%%%
\subsection{A Simple Optimal-Control Problem}
\label{subsec:simple-optimal-control-problem}
%%%%%%%%%%%%%%%%%%%%%%%%%%%%%%%%%%%%%%%%%%%%%%%%%%%%%%%%%%%%

Consider

\[
\boxed{
\dot{x}=u,
}
\]

with

\[
x(0)=x_0.
\]

Suppose the objective is

\[
\boxed{
J
=
\frac12
x(T)^2
+
\frac12
\int_0^T
u(t)^2\,dt.
}
\]

The control balances:

\[
\boxed{
\text{terminal state error}
}
\]

against

\[
\boxed{
\text{control effort}.
}
\]

%%%%%%%%%%%%%%%%%%%%%%%%%%%%%%%%%%%%%%%%%%%%%%%%%%%%%%%%%%%%
\subsection{Optimal Control and Calculus of Variations}
\label{subsec:optimal-control-calculus-variations}
%%%%%%%%%%%%%%%%%%%%%%%%%%%%%%%%%%%%%%%%%%%%%%%%%%%%%%%%%%%%

Calculus of variations optimizes functionals such as

\[
\boxed{
J[y]
=
\int
F
(
t,
y,
\dot y
)
\,dt.
}
\]

Optimal control extends this framework by introducing:

\[
\boxed{
\text{state equations}
}
\]

and explicit

\[
\boxed{
\text{control variables}.
}
\]

The state dynamics act as differential constraints.

%%%%%%%%%%%%%%%%%%%%%%%%%%%%%%%%%%%%%%%%%%%%%%%%%%%%%%%%%%%%
\subsection{Hamiltonian}
\label{subsec:optimal-control-hamiltonian}
%%%%%%%%%%%%%%%%%%%%%%%%%%%%%%%%%%%%%%%%%%%%%%%%%%%%%%%%%%%%

For the minimization convention used in this section, define the
Hamiltonian

\[
\boxed{
H(x,u,\lambda,t)
=
L(x,u,t)
+
\lambda^T
f(x,u,t),
}
\]

where

\[
\boxed{
\lambda(t)
\in
\mathbb{R}^n
}
\]

is the costate or adjoint vector.

%%%%%%%%%%%%%%%%%%%%%%%%%%%%%%%%%%%%%%%%%%%%%%%%%%%%%%%%%%%%
\subsection{Costate Variables}
\label{subsec:costate-variables}
%%%%%%%%%%%%%%%%%%%%%%%%%%%%%%%%%%%%%%%%%%%%%%%%%%%%%%%%%%%%

The costate variable acts as a time-dependent multiplier for the state
equation.

It measures, in a local sensitivity sense, the marginal value of changing
the state along an optimal trajectory.

Costates are analogous to Lagrange multipliers in finite-dimensional
optimization.

%%%%%%%%%%%%%%%%%%%%%%%%%%%%%%%%%%%%%%%%%%%%%%%%%%%%%%%%%%%%
\subsection{Pontryagin's Principle: Minimization Convention}
\label{subsec:pontryagin-principle}
%%%%%%%%%%%%%%%%%%%%%%%%%%%%%%%%%%%%%%%%%%%%%%%%%%%%%%%%%%%%

Consider

\[
\min
\left[
\Phi(x(T))
+
\int_{t_0}^{T}
L(x,u,t)\,dt
\right]
\]

subject to

\[
\dot{x}=f(x,u,t).
\]

Under appropriate regularity and optimality assumptions, an optimal
trajectory satisfies necessary conditions involving

\[
H
=
L+\lambda^Tf.
\]

The principal equations are

\[
\boxed{
\dot{x}
=
\frac{\partial H}{\partial\lambda}
=
f(x,u,t),
}
\]

\[
\boxed{
\dot{\lambda}
=
-
\frac{\partial H}{\partial x},
}
\]

and, for an interior control,

\[
\boxed{
\frac{\partial H}{\partial u}
=
0.
}
\]

More generally,

\[
\boxed{
u^*(t)
\in
\operatorname*{arg\,min}_{u\in U}
H
(
x^*(t),
u,
\lambda^*(t),
t
).
}
\]

%%%%%%%%%%%%%%%%%%%%%%%%%%%%%%%%%%%%%%%%%%%%%%%%%%%%%%%%%%%%
\subsection{Sign Conventions}
\label{subsec:optimal-control-sign-conventions}
%%%%%%%%%%%%%%%%%%%%%%%%%%%%%%%%%%%%%%%%%%%%%%%%%%%%%%%%%%%%

Different texts define the Hamiltonian differently.

This section uses

\[
\boxed{
H
=
L+\lambda^Tf
}
\]

for a minimization problem and therefore minimizes \(H\) with respect to
the control.

Another common convention uses

\[
H
=
-L+\lambda^Tf
\]

and maximizes the Hamiltonian.

Both conventions are valid when used consistently.

%%%%%%%%%%%%%%%%%%%%%%%%%%%%%%%%%%%%%%%%%%%%%%%%%%%%%%%%%%%%
\subsection{State Equation From the Hamiltonian}
\label{subsec:hamiltonian-state-equation}
%%%%%%%%%%%%%%%%%%%%%%%%%%%%%%%%%%%%%%%%%%%%%%%%%%%%%%%%%%%%

Because

\[
H
=
L+\lambda^Tf,
\]

we have

\[
\frac{\partial H}{\partial\lambda}
=
f.
\]

Thus the state equation is recovered as

\[
\boxed{
\dot{x}
=
H_{\lambda}.
}
\]

%%%%%%%%%%%%%%%%%%%%%%%%%%%%%%%%%%%%%%%%%%%%%%%%%%%%%%%%%%%%
\subsection{Costate Equation}
\label{subsec:costate-equation}
%%%%%%%%%%%%%%%%%%%%%%%%%%%%%%%%%%%%%%%%%%%%%%%%%%%%%%%%%%%%

The costate equation is

\[
\boxed{
\dot{\lambda}
=
-
H_x.
}
\]

In components,

\[
\boxed{
\dot{\lambda}_i
=
-
\frac{\partial H}{\partial x_i}.
}
\]

This differential equation is typically solved backward from a terminal
condition.

%%%%%%%%%%%%%%%%%%%%%%%%%%%%%%%%%%%%%%%%%%%%%%%%%%%%%%%%%%%%
\subsection{Terminal Transversality Condition}
\label{subsec:terminal-transversality}
%%%%%%%%%%%%%%%%%%%%%%%%%%%%%%%%%%%%%%%%%%%%%%%%%%%%%%%%%%%%

If terminal state \(x(T)\) is free and terminal time \(T\) is fixed, then

\[
\boxed{
\lambda(T)
=
\nabla_x
\Phi
(
x(T)
).
}
\]

If the terminal state is fixed, this costate terminal condition is replaced
by the corresponding endpoint restriction and multiplier structure.

%%%%%%%%%%%%%%%%%%%%%%%%%%%%%%%%%%%%%%%%%%%%%%%%%%%%%%%%%%%%
\subsection{Interior Control Condition}
\label{subsec:interior-control-condition}
%%%%%%%%%%%%%%%%%%%%%%%%%%%%%%%%%%%%%%%%%%%%%%%%%%%%%%%%%%%%

If the optimal control lies in the interior of the admissible set and the
Hamiltonian is differentiable with respect to \(u\), then

\[
\boxed{
H_u
=
0.
}
\]

This is analogous to a stationarity condition in ordinary optimization.

%%%%%%%%%%%%%%%%%%%%%%%%%%%%%%%%%%%%%%%%%%%%%%%%%%%%%%%%%%%%
\subsection{Control Constraints}
\label{subsec:optimal-control-control-constraints}
%%%%%%%%%%%%%%%%%%%%%%%%%%%%%%%%%%%%%%%%%%%%%%%%%%%%%%%%%%%%

Suppose

\[
\boxed{
u_{\min}
\leq
u(t)
\leq
u_{\max}.
}
\]

Then the stationarity equation

\[
H_u=0
\]

may produce an inadmissible value.

The correct condition is

\[
\boxed{
u^*(t)
\in
\operatorname*{arg\,min}_{
u_{\min}\leq u\leq u_{\max}
}
H.
}
\]

%%%%%%%%%%%%%%%%%%%%%%%%%%%%%%%%%%%%%%%%%%%%%%%%%%%%%%%%%%%%
\subsection{Worked Example: Quadratic Control Effort}
\label{subsec:worked-quadratic-control}
%%%%%%%%%%%%%%%%%%%%%%%%%%%%%%%%%%%%%%%%%%%%%%%%%%%%%%%%%%%%

\begin{workedexamplebox}{Hamiltonian Stationarity}

Suppose

\[
\dot{x}=u,
\]

and

\[
J
=
\frac12x(T)^2
+
\frac12
\int_0^T
u^2\,dt.
\]

The Hamiltonian is

\[
H
=
\frac12u^2
+
\lambda u.
\]

The stationarity condition is

\[
\frac{\partial H}{\partial u}
=
u+\lambda
=
0.
\]

Thus,

\begin{answerbox}
\[
\boxed{
u^*
=
-\lambda.
}
\]
\end{answerbox}

Because \(H\) does not depend explicitly on \(x\),

\[
\dot{\lambda}=0.
\]

Therefore the costate is constant.

\end{workedexamplebox}

%%%%%%%%%%%%%%%%%%%%%%%%%%%%%%%%%%%%%%%%%%%%%%%%%%%%%%%%%%%%
\subsection{Solving the Previous Example}
\label{subsec:solving-quadratic-control-example}
%%%%%%%%%%%%%%%%%%%%%%%%%%%%%%%%%%%%%%%%%%%%%%%%%%%%%%%%%%%%

Since

\[
\dot{\lambda}=0,
\]

write

\[
\lambda(t)=c.
\]

Then

\[
u^*=-c.
\]

The state satisfies

\[
\dot{x}=-c,
\]

so

\[
x(t)
=
x_0-ct.
\]

The terminal condition is

\[
\lambda(T)
=
x(T).
\]

Therefore,

\[
c
=
x_0-cT.
\]

Thus,

\[
c(1+T)
=
x_0,
\]

so

\[
\boxed{
c
=
\frac{x_0}{1+T}.
}
\]

Hence,

\[
\boxed{
u^*(t)
=
-
\frac{x_0}{1+T}.
}
\]

%%%%%%%%%%%%%%%%%%%%%%%%%%%%%%%%%%%%%%%%%%%%%%%%%%%%%%%%%%%%
\subsection{Optimal State Trajectory}
\label{subsec:optimal-state-trajectory-example}
%%%%%%%%%%%%%%%%%%%%%%%%%%%%%%%%%%%%%%%%%%%%%%%%%%%%%%%%%%%%

Substituting the optimal control,

\[
\dot{x}
=
-
\frac{x_0}{1+T}.
\]

Therefore,

\[
\boxed{
x^*(t)
=
x_0
-
\frac{x_0}{1+T}t.
}
\]

At the final time,

\[
\boxed{
x^*(T)
=
\frac{x_0}{1+T}.
}
\]

The terminal state is not forced exactly to zero because doing so would
require additional control effort.

%%%%%%%%%%%%%%%%%%%%%%%%%%%%%%%%%%%%%%%%%%%%%%%%%%%%%%%%%%%%
\subsection{Two-Point Boundary-Value Problem}
\label{subsec:two-point-boundary-value}
%%%%%%%%%%%%%%%%%%%%%%%%%%%%%%%%%%%%%%%%%%%%%%%%%%%%%%%%%%%%

Pontryagin conditions frequently create differential equations with
boundary conditions at both ends of the time interval:

\[
\boxed{
x(t_0)=x_0,
}
\]

and

\[
\boxed{
\lambda(T)
=
\nabla\Phi(x(T)).
}
\]

Such systems are called two-point boundary-value problems.

They are generally more difficult than ordinary initial-value problems.

%%%%%%%%%%%%%%%%%%%%%%%%%%%%%%%%%%%%%%%%%%%%%%%%%%%%%%%%%%%%
\subsection{Hamiltonian Conservation}
\label{subsec:hamiltonian-conservation}
%%%%%%%%%%%%%%%%%%%%%%%%%%%%%%%%%%%%%%%%%%%%%%%%%%%%%%%%%%%%

For sufficiently smooth autonomous optimal-control problems, where the
Hamiltonian has no explicit time dependence,

\[
\boxed{
\frac{\partial H}{\partial t}=0,
}
\]

the Hamiltonian is constant along suitable unconstrained extremals.

This property depends on the precise optimal-control setting and
regularity assumptions.

%%%%%%%%%%%%%%%%%%%%%%%%%%%%%%%%%%%%%%%%%%%%%%%%%%%%%%%%%%%%
\subsection{Free Final Time}
\label{subsec:free-final-time}
%%%%%%%%%%%%%%%%%%%%%%%%%%%%%%%%%%%%%%%%%%%%%%%%%%%%%%%%%%%%

If terminal time \(T\) is free, an additional transversality condition
appears.

For the convention

\[
H=L+\lambda^Tf
\]

and terminal cost

\[
\Phi(x(T),T),
\]

a standard condition is

\[
\boxed{
H(T)
+
\frac{\partial\Phi}{\partial T}
=
0,
}
\]

subject to the assumptions of the specific endpoint problem.

%%%%%%%%%%%%%%%%%%%%%%%%%%%%%%%%%%%%%%%%%%%%%%%%%%%%%%%%%%%%
\subsection{Minimum-Time Control}
\label{subsec:minimum-time-control}
%%%%%%%%%%%%%%%%%%%%%%%%%%%%%%%%%%%%%%%%%%%%%%%%%%%%%%%%%%%%

A minimum-time problem seeks

\[
\boxed{
\min T
}
\]

subject to dynamical constraints and a target condition.

It can be written with running cost

\[
\boxed{
L=1,
}
\]

so that

\[
\boxed{
J
=
\int_0^T1\,dt
=
T.
}
\]

Control bounds are often essential in minimum-time problems.

%%%%%%%%%%%%%%%%%%%%%%%%%%%%%%%%%%%%%%%%%%%%%%%%%%%%%%%%%%%%
\subsection{Bang-Bang Control}
\label{subsec:bang-bang-control}
%%%%%%%%%%%%%%%%%%%%%%%%%%%%%%%%%%%%%%%%%%%%%%%%%%%%%%%%%%%%

Suppose the Hamiltonian is linear in a bounded scalar control:

\[
\boxed{
H
=
H_0(x,\lambda,t)
+
\psi(t)u,
}
\]

with

\[
u_{\min}
\leq
u
\leq
u_{\max}.
\]

To minimize \(H\),

\[
\boxed{
u^*
=
u_{\min}
}
\]

when

\[
\psi(t)>0,
\]

and

\[
\boxed{
u^*
=
u_{\max}
}
\]

when

\[
\psi(t)<0.
\]

Thus the control switches between its extreme values.

%%%%%%%%%%%%%%%%%%%%%%%%%%%%%%%%%%%%%%%%%%%%%%%%%%%%%%%%%%%%
\subsection{Switching Function}
\label{subsec:switching-function}
%%%%%%%%%%%%%%%%%%%%%%%%%%%%%%%%%%%%%%%%%%%%%%%%%%%%%%%%%%%%

The coefficient

\[
\boxed{
\psi(t)
=
\frac{\partial H}{\partial u}
}
\]

for a Hamiltonian affine in \(u\) is called a switching function.

Its sign determines which boundary control minimizes the Hamiltonian.

A switching time occurs when the sign changes.

%%%%%%%%%%%%%%%%%%%%%%%%%%%%%%%%%%%%%%%%%%%%%%%%%%%%%%%%%%%%
\subsection{Singular Arcs}
\label{subsec:singular-arcs}
%%%%%%%%%%%%%%%%%%%%%%%%%%%%%%%%%%%%%%%%%%%%%%%%%%%%%%%%%%%%

If

\[
\boxed{
\psi(t)=0
}
\]

over an interval, the first-order Hamiltonian condition does not determine
the control directly.

Such an interval may be a singular arc.

Higher derivatives of the switching condition and additional optimality
conditions may be required.

%%%%%%%%%%%%%%%%%%%%%%%%%%%%%%%%%%%%%%%%%%%%%%%%%%%%%%%%%%%%
\subsection{Worked Example: Bang-Bang Decision}
\label{subsec:worked-bang-bang}
%%%%%%%%%%%%%%%%%%%%%%%%%%%%%%%%%%%%%%%%%%%%%%%%%%%%%%%%%%%%

\begin{workedexamplebox}{Bounded Linear Hamiltonian}

Suppose

\[
H
=
5
+
3u,
\]

with

\[
-2
\leq
u
\leq
2.
\]

Because the coefficient of \(u\) is positive, minimizing \(H\) requires
the smallest admissible control.

Therefore,

\begin{answerbox}
\[
\boxed{
u^*=-2.
}
\]
\end{answerbox}

If instead the coefficient were negative, the minimizing control would be
the upper bound.

\end{workedexamplebox}

%%%%%%%%%%%%%%%%%%%%%%%%%%%%%%%%%%%%%%%%%%%%%%%%%%%%%%%%%%%%
\subsection{Minimum-Energy Control}
\label{subsec:minimum-energy-control}
%%%%%%%%%%%%%%%%%%%%%%%%%%%%%%%%%%%%%%%%%%%%%%%%%%%%%%%%%%%%

A common objective is

\[
\boxed{
J
=
\frac12
\int_0^T
u(t)^Tu(t)
\,dt.
}
\]

This minimizes control energy or effort.

Such objectives strongly favor smooth moderate controls rather than
bang-bang inputs.

%%%%%%%%%%%%%%%%%%%%%%%%%%%%%%%%%%%%%%%%%%%%%%%%%%%%%%%%%%%%
\subsection{Tracking Problems}
\label{subsec:optimal-control-tracking}
%%%%%%%%%%%%%%%%%%%%%%%%%%%%%%%%%%%%%%%%%%%%%%%%%%%%%%%%%%%%

Suppose a desired state trajectory is

\[
x_r(t).
\]

A quadratic tracking objective is

\[
\boxed{
J
=
\frac12
(
x(T)-x_r(T)
)^T
F
(
x(T)-x_r(T)
)
}
\]

\[
\boxed{
\phantom{J={}}
+
\frac12
\int_0^T
\left[
(
x-x_r
)^T
Q
(
x-x_r
)
+
u^TRu
\right]dt.
}
\]

This balances tracking accuracy and control effort.

%%%%%%%%%%%%%%%%%%%%%%%%%%%%%%%%%%%%%%%%%%%%%%%%%%%%%%%%%%%%
\subsection{Linear-Quadratic Regulator}
\label{subsec:linear-quadratic-regulator}
%%%%%%%%%%%%%%%%%%%%%%%%%%%%%%%%%%%%%%%%%%%%%%%%%%%%%%%%%%%%

Consider the linear system

\[
\boxed{
\dot{x}
=
Ax+Bu,
}
\]

with quadratic objective

\[
\boxed{
J
=
\frac12x(T)^TFx(T)
+
\frac12
\int_0^T
\left(
x^TQx
+
u^TRu
\right)dt.
}
\]

Assume typically

\[
\boxed{
Q\succeq0,
\qquad
F\succeq0,
\qquad
R\succ0.
}
\]

This is the finite-horizon Linear-Quadratic Regulator, abbreviated LQR.

%%%%%%%%%%%%%%%%%%%%%%%%%%%%%%%%%%%%%%%%%%%%%%%%%%%%%%%%%%%%
\subsection{LQR Hamiltonian}
\label{subsec:lqr-hamiltonian}
%%%%%%%%%%%%%%%%%%%%%%%%%%%%%%%%%%%%%%%%%%%%%%%%%%%%%%%%%%%%

The Hamiltonian is

\[
\boxed{
H
=
\frac12x^TQx
+
\frac12u^TRu
+
\lambda^T
(
Ax+Bu
).
}
\]

Stationarity with respect to \(u\) gives

\[
Ru
+
B^T\lambda
=
0.
\]

Therefore,

\[
\boxed{
u^*
=
-
R^{-1}B^T\lambda.
}
\]

%%%%%%%%%%%%%%%%%%%%%%%%%%%%%%%%%%%%%%%%%%%%%%%%%%%%%%%%%%%%
\subsection{LQR Costate Equation}
\label{subsec:lqr-costate}
%%%%%%%%%%%%%%%%%%%%%%%%%%%%%%%%%%%%%%%%%%%%%%%%%%%%%%%%%%%%

The costate equation is

\[
\dot{\lambda}
=
-
Qx
-
A^T\lambda.
\]

The terminal condition is

\[
\boxed{
\lambda(T)
=
Fx(T).
}
\]

This suggests a relationship of the form

\[
\boxed{
\lambda(t)
=
P(t)x(t).
}
\]

%%%%%%%%%%%%%%%%%%%%%%%%%%%%%%%%%%%%%%%%%%%%%%%%%%%%%%%%%%%%
\subsection{Riccati Differential Equation}
\label{subsec:riccati-differential-equation}
%%%%%%%%%%%%%%%%%%%%%%%%%%%%%%%%%%%%%%%%%%%%%%%%%%%%%%%%%%%%

Substituting

\[
\lambda=Px
\]

into the LQR conditions yields

\[
\boxed{
-\dot{P}
=
A^TP
+
PA
-
PBR^{-1}B^TP
+
Q,
}
\]

with terminal condition

\[
\boxed{
P(T)=F.
}
\]

This is the continuous-time Riccati differential equation.

%%%%%%%%%%%%%%%%%%%%%%%%%%%%%%%%%%%%%%%%%%%%%%%%%%%%%%%%%%%%
\subsection{LQR Feedback Law}
\label{subsec:lqr-feedback-law}
%%%%%%%%%%%%%%%%%%%%%%%%%%%%%%%%%%%%%%%%%%%%%%%%%%%%%%%%%%%%

The optimal control becomes

\[
\boxed{
u^*(t)
=
-
R^{-1}
B^T
P(t)
x(t).
}
\]

Define

\[
\boxed{
K(t)
=
R^{-1}B^TP(t).
}
\]

Then

\[
\boxed{
u^*(t)
=
-K(t)x(t).
}
\]

Thus finite-horizon LQR naturally produces linear state feedback.

%%%%%%%%%%%%%%%%%%%%%%%%%%%%%%%%%%%%%%%%%%%%%%%%%%%%%%%%%%%%
\subsection{Interpretation of LQR Weights}
\label{subsec:lqr-weight-interpretation}
%%%%%%%%%%%%%%%%%%%%%%%%%%%%%%%%%%%%%%%%%%%%%%%%%%%%%%%%%%%%

The matrix

\[
Q
\]

penalizes state deviation.

The matrix

\[
R
\]

penalizes control effort.

Larger state penalties encourage more aggressive correction.

Larger control penalties discourage large inputs.

Thus LQR explicitly models a performance-effort tradeoff.

%%%%%%%%%%%%%%%%%%%%%%%%%%%%%%%%%%%%%%%%%%%%%%%%%%%%%%%%%%%%
\subsection{Scalar LQR Example}
\label{subsec:scalar-lqr-example}
%%%%%%%%%%%%%%%%%%%%%%%%%%%%%%%%%%%%%%%%%%%%%%%%%%%%%%%%%%%%

Consider

\[
\dot{x}
=
ax+bu,
\]

and

\[
J
=
\frac12Fx(T)^2
+
\frac12
\int_0^T
\left(
qx^2+ru^2
\right)dt.
\]

The Riccati equation becomes

\[
\boxed{
-\dot{P}
=
2aP
-
\frac{b^2}{r}P^2
+
q,
}
\]

with

\[
\boxed{
P(T)=F.
}
\]

The feedback law is

\[
\boxed{
u^*
=
-
\frac{b}{r}
Px.
}
\]

%%%%%%%%%%%%%%%%%%%%%%%%%%%%%%%%%%%%%%%%%%%%%%%%%%%%%%%%%%%%
\subsection{Infinite-Horizon LQR}
\label{subsec:infinite-horizon-lqr}
%%%%%%%%%%%%%%%%%%%%%%%%%%%%%%%%%%%%%%%%%%%%%%%%%%%%%%%%%%%%

Consider

\[
\boxed{
J
=
\frac12
\int_0^\infty
\left(
x^TQx
+
u^TRu
\right)dt.
}
\]

Under appropriate stabilizability and detectability assumptions, the
steady-state matrix \(P\) satisfies the algebraic Riccati equation

\[
\boxed{
A^TP
+
PA
-
PBR^{-1}B^TP
+
Q
=
0.
}
\]

%%%%%%%%%%%%%%%%%%%%%%%%%%%%%%%%%%%%%%%%%%%%%%%%%%%%%%%%%%%%
\subsection{Infinite-Horizon Feedback}
\label{subsec:infinite-horizon-feedback}
%%%%%%%%%%%%%%%%%%%%%%%%%%%%%%%%%%%%%%%%%%%%%%%%%%%%%%%%%%%%

The resulting feedback law is

\[
\boxed{
u^*
=
-Kx,
}
\]

where

\[
\boxed{
K
=
R^{-1}B^TP.
}
\]

Unlike finite-horizon LQR, \(K\) is constant when the system and objective
are time invariant.

%%%%%%%%%%%%%%%%%%%%%%%%%%%%%%%%%%%%%%%%%%%%%%%%%%%%%%%%%%%%
\subsection{Controllability}
\label{subsec:optimal-control-controllability}
%%%%%%%%%%%%%%%%%%%%%%%%%%%%%%%%%%%%%%%%%%%%%%%%%%%%%%%%%%%%

Before optimizing a control problem, one must determine whether the desired
state can be reached at all.

For the linear system

\[
\dot{x}=Ax+Bu,
\]

define the controllability matrix

\[
\boxed{
\mathcal{C}
=
\begin{bmatrix}
B
&
AB
&
A^2B
&
\cdots
&
A^{n-1}B
\end{bmatrix}.
}
\]

The pair \((A,B)\) is controllable if

\[
\boxed{
\operatorname{rank}
(
\mathcal{C}
)
=
n.
}
\]

%%%%%%%%%%%%%%%%%%%%%%%%%%%%%%%%%%%%%%%%%%%%%%%%%%%%%%%%%%%%
\subsection{Feasibility Before Optimality}
\label{subsec:feasibility-before-optimality}
%%%%%%%%%%%%%%%%%%%%%%%%%%%%%%%%%%%%%%%%%%%%%%%%%%%%%%%%%%%%

An optimal-control objective cannot overcome impossible dynamics.

Thus the logical order is:

\[
\boxed{
\text{Can the system reach the required state?}
}
\]

before asking:

\[
\boxed{
\text{What is the best way to reach it?}
}
\]

Controllability is therefore fundamentally a feasibility concept.

%%%%%%%%%%%%%%%%%%%%%%%%%%%%%%%%%%%%%%%%%%%%%%%%%%%%%%%%%%%%
\subsection{Discrete-Time Optimal Control}
\label{subsec:discrete-time-optimal-control}
%%%%%%%%%%%%%%%%%%%%%%%%%%%%%%%%%%%%%%%%%%%%%%%%%%%%%%%%%%%%

A discrete-time system has dynamics

\[
\boxed{
x_{k+1}
=
f_k
(
x_k,
u_k
).
}
\]

A finite-horizon objective may be

\[
\boxed{
J
=
\Phi(x_N)
+
\sum_{k=0}^{N-1}
L_k
(
x_k,
u_k
).
}
\]

This form connects directly with dynamic programming.

%%%%%%%%%%%%%%%%%%%%%%%%%%%%%%%%%%%%%%%%%%%%%%%%%%%%%%%%%%%%
\subsection{Bellman Principle}
\label{subsec:optimal-control-bellman-principle}
%%%%%%%%%%%%%%%%%%%%%%%%%%%%%%%%%%%%%%%%%%%%%%%%%%%%%%%%%%%%

Bellman's principle of optimality states that the remaining decisions of an
optimal policy must themselves be optimal for the state reached after the
first decision.

This produces recursive optimization.

%%%%%%%%%%%%%%%%%%%%%%%%%%%%%%%%%%%%%%%%%%%%%%%%%%%%%%%%%%%%
\subsection{Discrete-Time Value Function}
\label{subsec:discrete-value-function}
%%%%%%%%%%%%%%%%%%%%%%%%%%%%%%%%%%%%%%%%%%%%%%%%%%%%%%%%%%%%

Define

\[
\boxed{
V_k(x)
=
\text{minimum future cost starting from state }x\text{ at stage }k.
}
\]

Then

\[
\boxed{
V_N(x)
=
\Phi(x).
}
\]

The Bellman recursion is

\[
\boxed{
V_k(x)
=
\min_{u\in U_k(x)}
\left[
L_k(x,u)
+
V_{k+1}
(
f_k(x,u)
)
\right].
}
\]

%%%%%%%%%%%%%%%%%%%%%%%%%%%%%%%%%%%%%%%%%%%%%%%%%%%%%%%%%%%%
\subsection{Feedback Policy From Bellman Recursion}
\label{subsec:feedback-from-bellman}
%%%%%%%%%%%%%%%%%%%%%%%%%%%%%%%%%%%%%%%%%%%%%%%%%%%%%%%%%%%%

The minimizing control is

\[
\boxed{
u_k^*(x)
\in
\operatorname*{arg\,min}_{u}
\left[
L_k(x,u)
+
V_{k+1}
(
f_k(x,u)
)
\right].
}
\]

Thus dynamic programming gives a feedback policy directly:

\[
\boxed{
u_k^*
=
\pi_k(x_k).
}
\]

%%%%%%%%%%%%%%%%%%%%%%%%%%%%%%%%%%%%%%%%%%%%%%%%%%%%%%%%%%%%
\subsection{Continuous-Time Value Function}
\label{subsec:continuous-value-function}
%%%%%%%%%%%%%%%%%%%%%%%%%%%%%%%%%%%%%%%%%%%%%%%%%%%%%%%%%%%%

For continuous time, define

\[
\boxed{
V(x,t)
=
\inf_{u(\cdot)}
\left[
\Phi(x(T))
+
\int_t^T
L
(
x(s),
u(s),
s
)
\,ds
\right],
}
\]

where the trajectory starts from state \(x\) at time \(t\).

%%%%%%%%%%%%%%%%%%%%%%%%%%%%%%%%%%%%%%%%%%%%%%%%%%%%%%%%%%%%
\subsection{Hamilton--Jacobi--Bellman Equation}
\label{subsec:hjb-equation}
%%%%%%%%%%%%%%%%%%%%%%%%%%%%%%%%%%%%%%%%%%%%%%%%%%%%%%%%%%%%

Under appropriate smoothness assumptions, the value function satisfies the
Hamilton--Jacobi--Bellman equation

\[
\boxed{
-
V_t(x,t)
=
\min_{u\in U}
\left\{
L(x,u,t)
+
\nabla_xV(x,t)^T
f(x,u,t)
\right\}.
}
\]

The terminal condition is

\[
\boxed{
V(x,T)
=
\Phi(x).
}
\]

%%%%%%%%%%%%%%%%%%%%%%%%%%%%%%%%%%%%%%%%%%%%%%%%%%%%%%%%%%%%
\subsection{HJB as an Optimization Equation}
\label{subsec:hjb-optimization-equation}
%%%%%%%%%%%%%%%%%%%%%%%%%%%%%%%%%%%%%%%%%%%%%%%%%%%%%%%%%%%%

The expression

\[
L
+
\nabla V^Tf
\]

has the same structure as the Hamiltonian

\[
H
=
L+\lambda^Tf.
\]

Indeed, when the value function is sufficiently smooth,

\[
\boxed{
\lambda
=
\nabla_xV
}
\]

along an optimal trajectory under standard conditions.

Thus PMP and HJB are closely related.

%%%%%%%%%%%%%%%%%%%%%%%%%%%%%%%%%%%%%%%%%%%%%%%%%%%%%%%%%%%%
\subsection{Optimal Feedback From HJB}
\label{subsec:optimal-feedback-hjb}
%%%%%%%%%%%%%%%%%%%%%%%%%%%%%%%%%%%%%%%%%%%%%%%%%%%%%%%%%%%%

If the minimizer is unique,

\[
\boxed{
u^*(x,t)
=
\operatorname*{arg\,min}_{u\in U}
\left\{
L(x,u,t)
+
\nabla V(x,t)^Tf(x,u,t)
\right\}.
}
\]

Therefore HJB produces an optimal control as a function of the state.

%%%%%%%%%%%%%%%%%%%%%%%%%%%%%%%%%%%%%%%%%%%%%%%%%%%%%%%%%%%%
\subsection{Worked Example: HJB for an Integrator}
\label{subsec:worked-hjb-integrator}
%%%%%%%%%%%%%%%%%%%%%%%%%%%%%%%%%%%%%%%%%%%%%%%%%%%%%%%%%%%%

\begin{workedexamplebox}{Quadratic Control of an Integrator}

Consider

\[
\dot{x}=u,
\]

with

\[
J
=
\frac12x(T)^2
+
\frac12
\int_t^T
u^2\,ds.
\]

The HJB equation is

\[
-V_t
=
\min_u
\left\{
\frac12u^2
+
V_xu
\right\}.
\]

Differentiate with respect to \(u\):

\[
u+V_x=0.
\]

Therefore,

\begin{answerbox}
\[
\boxed{
u^*
=
-V_x.
}
\]
\end{answerbox}

Substitution gives

\[
\boxed{
-V_t
=
-\frac12V_x^2.
}
\]

Equivalently,

\[
\boxed{
V_t
=
\frac12V_x^2.
}
\]

\end{workedexamplebox}

%%%%%%%%%%%%%%%%%%%%%%%%%%%%%%%%%%%%%%%%%%%%%%%%%%%%%%%%%%%%
\subsection{PMP Versus HJB}
\label{subsec:pmp-versus-hjb}
%%%%%%%%%%%%%%%%%%%%%%%%%%%%%%%%%%%%%%%%%%%%%%%%%%%%%%%%%%%%

Pontryagin's principle gives necessary conditions along candidate optimal
trajectories.

The HJB equation characterizes the optimal value function under suitable
conditions and produces state-feedback policies.

Thus:

\[
\boxed{
\text{PMP}
\rightarrow
\text{trajectory-based necessary conditions},
}
\]

while

\[
\boxed{
\text{HJB}
\rightarrow
\text{value-function dynamic programming}.
}
\]

%%%%%%%%%%%%%%%%%%%%%%%%%%%%%%%%%%%%%%%%%%%%%%%%%%%%%%%%%%%%
\subsection{Curse of Dimensionality}
\label{subsec:optimal-control-curse-dimensionality}
%%%%%%%%%%%%%%%%%%%%%%%%%%%%%%%%%%%%%%%%%%%%%%%%%%%%%%%%%%%%

If the state has dimension \(n\), numerically approximating

\[
V(x,t)
\]

on a grid can require exponentially many state points.

This is the curse of dimensionality.

It makes classical HJB methods difficult for high-dimensional systems.

%%%%%%%%%%%%%%%%%%%%%%%%%%%%%%%%%%%%%%%%%%%%%%%%%%%%%%%%%%%%
\subsection{State Constraints}
\label{subsec:optimal-control-state-constraints}
%%%%%%%%%%%%%%%%%%%%%%%%%%%%%%%%%%%%%%%%%%%%%%%%%%%%%%%%%%%%

A state may be required to satisfy

\[
\boxed{
g(x(t),t)
\leq
0.
}
\]

Examples include:

\[
\boxed{
\text{altitude limits},
}
\]

\[
\boxed{
\text{temperature limits},
}
\]

\[
\boxed{
\text{inventory bounds},
}
\]

or

\[
\boxed{
\text{safety regions}.
}
\]

State constraints introduce additional multiplier and boundary conditions.

%%%%%%%%%%%%%%%%%%%%%%%%%%%%%%%%%%%%%%%%%%%%%%%%%%%%%%%%%%%%
\subsection{Path Constraints}
\label{subsec:path-constraints}
%%%%%%%%%%%%%%%%%%%%%%%%%%%%%%%%%%%%%%%%%%%%%%%%%%%%%%%%%%%%

A general path constraint has form

\[
\boxed{
g
(
x(t),
u(t),
t
)
\leq
0.
}
\]

Unlike endpoint constraints, it must hold throughout the entire time
interval.

Examples include actuator limits and safety constraints.

%%%%%%%%%%%%%%%%%%%%%%%%%%%%%%%%%%%%%%%%%%%%%%%%%%%%%%%%%%%%
\subsection{Endpoint Constraints}
\label{subsec:endpoint-constraints}
%%%%%%%%%%%%%%%%%%%%%%%%%%%%%%%%%%%%%%%%%%%%%%%%%%%%%%%%%%%%

An endpoint constraint may have form

\[
\boxed{
\Psi
(
x(t_0),
x(T),
T
)
=
0.
}
\]

Endpoint multipliers then modify the transversality conditions.

Optimal-control boundary conditions must therefore be derived from the
specific endpoint structure.

%%%%%%%%%%%%%%%%%%%%%%%%%%%%%%%%%%%%%%%%%%%%%%%%%%%%%%%%%%%%
\subsection{Path-Constrained Hamiltonian}
\label{subsec:path-constrained-hamiltonian}
%%%%%%%%%%%%%%%%%%%%%%%%%%%%%%%%%%%%%%%%%%%%%%%%%%%%%%%%%%%%

For inequality path constraint

\[
g(x,u,t)\leq0,
\]

one may introduce multiplier

\[
\mu(t)\geq0.
\]

An augmented Hamiltonian has form

\[
\boxed{
\mathcal{H}
=
L
+
\lambda^Tf
+
\mu^Tg.
}
\]

Complementarity conditions resemble finite-dimensional KKT conditions.

%%%%%%%%%%%%%%%%%%%%%%%%%%%%%%%%%%%%%%%%%%%%%%%%%%%%%%%%%%%%
\subsection{Complementarity Along a Trajectory}
\label{subsec:path-complementarity}
%%%%%%%%%%%%%%%%%%%%%%%%%%%%%%%%%%%%%%%%%%%%%%%%%%%%%%%%%%%%

For suitable path inequality constraints,

\[
\boxed{
\mu_i(t)\geq0,
}
\]

\[
\boxed{
g_i(x,u,t)\leq0,
}
\]

and

\[
\boxed{
\mu_i(t)
g_i(x,u,t)
=
0.
}
\]

Thus inactive constraints have zero multipliers, while active constraints
may influence the optimal control.

%%%%%%%%%%%%%%%%%%%%%%%%%%%%%%%%%%%%%%%%%%%%%%%%%%%%%%%%%%%%
\subsection{Terminal-State Penalties}
\label{subsec:terminal-state-penalties}
%%%%%%%%%%%%%%%%%%%%%%%%%%%%%%%%%%%%%%%%%%%%%%%%%%%%%%%%%%%%

Instead of requiring

\[
x(T)=x_T
\]

exactly, one may penalize terminal error:

\[
\boxed{
\Phi(x(T))
=
\frac12
(
x(T)-x_T
)^T
F
(
x(T)-x_T
).
}
\]

This permits a tradeoff between final accuracy and control effort.

%%%%%%%%%%%%%%%%%%%%%%%%%%%%%%%%%%%%%%%%%%%%%%%%%%%%%%%%%%%%
\subsection{Hard Versus Soft Constraints}
\label{subsec:hard-versus-soft-control}
%%%%%%%%%%%%%%%%%%%%%%%%%%%%%%%%%%%%%%%%%%%%%%%%%%%%%%%%%%%%

A hard constraint must be satisfied:

\[
\boxed{
g(x,u,t)\leq0.
}
\]

A soft constraint can be violated at a penalty.

For example,

\[
\boxed{
J
=
J_0
+
\rho
\int
[g(x,u,t)]_+^2dt.
}
\]

Soft constraints can improve numerical feasibility but no longer guarantee
exact satisfaction.

%%%%%%%%%%%%%%%%%%%%%%%%%%%%%%%%%%%%%%%%%%%%%%%%%%%%%%%%%%%%
\subsection{Discounted Infinite-Horizon Control}
\label{subsec:discounted-infinite-horizon-control}
%%%%%%%%%%%%%%%%%%%%%%%%%%%%%%%%%%%%%%%%%%%%%%%%%%%%%%%%%%%%

In economics and other applications, future costs may be discounted:

\[
\boxed{
J
=
\int_0^\infty
e^{-\rho t}
L(x(t),u(t))
\,dt,
}
\]

where

\[
\boxed{
\rho>0
}
\]

is a discount rate.

Future costs receive less weight than immediate costs.

%%%%%%%%%%%%%%%%%%%%%%%%%%%%%%%%%%%%%%%%%%%%%%%%%%%%%%%%%%%%
\subsection{Optimal Control in Economics}
\label{subsec:optimal-control-economics}
%%%%%%%%%%%%%%%%%%%%%%%%%%%%%%%%%%%%%%%%%%%%%%%%%%%%%%%%%%%%

An economic state might be capital

\[
K(t),
\]

with dynamics

\[
\boxed{
\dot{K}
=
F(K)
-
C
-
\delta K,
}
\]

where

\[
C(t)
\]

is consumption.

An objective may maximize discounted utility:

\[
\boxed{
\max
\int_0^\infty
e^{-\rho t}
U(C(t))
\,dt.
}
\]

Optimal-control methods then determine the consumption path.

%%%%%%%%%%%%%%%%%%%%%%%%%%%%%%%%%%%%%%%%%%%%%%%%%%%%%%%%%%%%
\subsection{Current-Value Hamiltonian Preview}
\label{subsec:current-value-hamiltonian}
%%%%%%%%%%%%%%%%%%%%%%%%%%%%%%%%%%%%%%%%%%%%%%%%%%%%%%%%%%%%

For discounted economic problems, one may transform the costate and define
a current-value Hamiltonian.

This removes the explicit discount factor from some first-order
conditions.

Care is required because the current-value costate equation differs from
the ordinary present-value equation.

%%%%%%%%%%%%%%%%%%%%%%%%%%%%%%%%%%%%%%%%%%%%%%%%%%%%%%%%%%%%
\subsection{Optimal Control in Biological Systems}
\label{subsec:optimal-control-biological}
%%%%%%%%%%%%%%%%%%%%%%%%%%%%%%%%%%%%%%%%%%%%%%%%%%%%%%%%%%%%

Optimal control can model interventions in systems such as:

\[
\boxed{
\text{epidemics},
}
\]

\[
\boxed{
\text{harvesting},
}
\]

\[
\boxed{
\text{drug treatment},
}
\]

or

\[
\boxed{
\text{population management}.
}
\]

The state equations are often nonlinear differential equations.

Controls represent intervention intensity.

%%%%%%%%%%%%%%%%%%%%%%%%%%%%%%%%%%%%%%%%%%%%%%%%%%%%%%%%%%%%
\subsection{Epidemic-Control Example}
\label{subsec:epidemic-control-example}
%%%%%%%%%%%%%%%%%%%%%%%%%%%%%%%%%%%%%%%%%%%%%%%%%%%%%%%%%%%%

Suppose infection level is \(I(t)\) and intervention intensity is \(u(t)\).

A schematic objective could be

\[
\boxed{
J
=
\int_0^T
\left[
aI(t)^2
+
bu(t)^2
\right]dt.
}
\]

The first term penalizes infection.

The second penalizes intervention cost.

The state equations constrain how intervention changes epidemic dynamics.

%%%%%%%%%%%%%%%%%%%%%%%%%%%%%%%%%%%%%%%%%%%%%%%%%%%%%%%%%%%%
\subsection{Trajectory Optimization}
\label{subsec:trajectory-optimization}
%%%%%%%%%%%%%%%%%%%%%%%%%%%%%%%%%%%%%%%%%%%%%%%%%%%%%%%%%%%%

Trajectory optimization is optimal control applied to motion planning.

The state may contain

\[
\boxed{
\text{position}
+
\text{velocity}
+
\text{orientation}.
}
\]

Controls may include forces, torques, steering, or thrust.

Constraints may enforce:

\[
\boxed{
\text{dynamics},
}
\]

\[
\boxed{
\text{obstacle avoidance},
}
\]

and

\[
\boxed{
\text{actuator limits}.
}
\]

%%%%%%%%%%%%%%%%%%%%%%%%%%%%%%%%%%%%%%%%%%%%%%%%%%%%%%%%%%%%
\subsection{Model Predictive Control}
\label{subsec:model-predictive-control}
%%%%%%%%%%%%%%%%%%%%%%%%%%%%%%%%%%%%%%%%%%%%%%%%%%%%%%%%%%%%

Model Predictive Control, abbreviated MPC, repeatedly solves a finite-
horizon optimal-control problem.

At each update:

\begin{enumerate}
    \item measure or estimate the current state;
    \item solve a finite-horizon control problem;
    \item apply only the first control action;
    \item advance the system;
    \item measure the new state;
    \item solve again.
\end{enumerate}

This is also called receding-horizon control.

%%%%%%%%%%%%%%%%%%%%%%%%%%%%%%%%%%%%%%%%%%%%%%%%%%%%%%%%%%%%
\subsection{Why MPC Uses Feedback}
\label{subsec:why-mpc-feedback}
%%%%%%%%%%%%%%%%%%%%%%%%%%%%%%%%%%%%%%%%%%%%%%%%%%%%%%%%%%%%

Although each finite-horizon optimization produces an open-loop sequence,
MPC repeatedly resolves the problem using the current measured state.

Thus the overall implementation creates feedback.

It can respond to:

\[
\boxed{
\text{disturbances},
}
\]

\[
\boxed{
\text{forecast changes},
}
\]

and

\[
\boxed{
\text{model mismatch}.
}
\]

%%%%%%%%%%%%%%%%%%%%%%%%%%%%%%%%%%%%%%%%%%%%%%%%%%%%%%%%%%%%
\subsection{MPC Objective}
\label{subsec:mpc-objective}
%%%%%%%%%%%%%%%%%%%%%%%%%%%%%%%%%%%%%%%%%%%%%%%%%%%%%%%%%%%%

A discrete-time MPC objective may be

\[
\boxed{
\sum_{k=0}^{N-1}
\left[
(x_k-x_r)^TQ(x_k-x_r)
+
u_k^TRu_k
\right]
+
(x_N-x_r)^TP(x_N-x_r).
}
\]

Constraints are imposed over the prediction horizon.

Only \(u_0\) is implemented before optimization is repeated.

%%%%%%%%%%%%%%%%%%%%%%%%%%%%%%%%%%%%%%%%%%%%%%%%%%%%%%%%%%%%
\subsection{Constrained MPC}
\label{subsec:constrained-mpc}
%%%%%%%%%%%%%%%%%%%%%%%%%%%%%%%%%%%%%%%%%%%%%%%%%%%%%%%%%%%%

MPC is particularly useful because it can handle constraints such as

\[
\boxed{
u_{\min}\leq u_k\leq u_{\max},
}
\]

and

\[
\boxed{
x_{\min}\leq x_k\leq x_{\max}.
}
\]

For linear dynamics and quadratic costs with linear constraints, each MPC
step may become a quadratic program.

%%%%%%%%%%%%%%%%%%%%%%%%%%%%%%%%%%%%%%%%%%%%%%%%%%%%%%%%%%%%
\subsection{Indirect Numerical Methods}
\label{subsec:indirect-optimal-control-methods}
%%%%%%%%%%%%%%%%%%%%%%%%%%%%%%%%%%%%%%%%%%%%%%%%%%%%%%%%%%%%

Indirect methods first derive necessary optimality conditions.

One then solves:

\[
\boxed{
\text{state equations}
+
\text{costate equations}
+
\text{optimality conditions}
+
\text{boundary conditions}.
}
\]

Examples include shooting and multiple-shooting methods.

%%%%%%%%%%%%%%%%%%%%%%%%%%%%%%%%%%%%%%%%%%%%%%%%%%%%%%%%%%%%
\subsection{Shooting Method}
\label{subsec:shooting-method}
%%%%%%%%%%%%%%%%%%%%%%%%%%%%%%%%%%%%%%%%%%%%%%%%%%%%%%%%%%%%

In a shooting method, unknown initial costate values are guessed.

The state-costate differential equations are integrated.

The resulting terminal error is used to update the initial costate guess.

Thus a boundary-value problem is transformed into a nonlinear root-finding
problem.

%%%%%%%%%%%%%%%%%%%%%%%%%%%%%%%%%%%%%%%%%%%%%%%%%%%%%%%%%%%%
\subsection{Sensitivity of Shooting Methods}
\label{subsec:shooting-sensitivity}
%%%%%%%%%%%%%%%%%%%%%%%%%%%%%%%%%%%%%%%%%%%%%%%%%%%%%%%%%%%%

Shooting can be extremely sensitive to the initial costate guess.

Small changes in the initial adjoint can lead to large terminal-state
changes.

This is especially problematic for:

\[
\boxed{
\text{long horizons},
}
\]

\[
\boxed{
\text{unstable dynamics},
}
\]

or

\[
\boxed{
\text{strong nonlinearities}.
}
\]

%%%%%%%%%%%%%%%%%%%%%%%%%%%%%%%%%%%%%%%%%%%%%%%%%%%%%%%%%%%%
\subsection{Multiple Shooting}
\label{subsec:multiple-shooting}
%%%%%%%%%%%%%%%%%%%%%%%%%%%%%%%%%%%%%%%%%%%%%%%%%%%%%%%%%%%%

Multiple shooting divides the horizon into several intervals.

Separate initial states or costates are introduced for each interval.

Continuity constraints then connect neighboring segments.

This often improves numerical stability compared with single shooting.

%%%%%%%%%%%%%%%%%%%%%%%%%%%%%%%%%%%%%%%%%%%%%%%%%%%%%%%%%%%%
\subsection{Direct Numerical Methods}
\label{subsec:direct-optimal-control-methods}
%%%%%%%%%%%%%%%%%%%%%%%%%%%%%%%%%%%%%%%%%%%%%%%%%%%%%%%%%%%%

Direct methods discretize the control problem before deriving optimization
conditions.

The result is a finite-dimensional nonlinear program or quadratic program.

Thus:

\[
\boxed{
\text{optimal control problem}
\rightarrow
\text{nonlinear programming problem}.
}
\]

%%%%%%%%%%%%%%%%%%%%%%%%%%%%%%%%%%%%%%%%%%%%%%%%%%%%%%%%%%%%
\subsection{Control Parameterization}
\label{subsec:control-parameterization}
%%%%%%%%%%%%%%%%%%%%%%%%%%%%%%%%%%%%%%%%%%%%%%%%%%%%%%%%%%%%

Approximate the control by finitely many parameters.

For piecewise-constant control,

\[
\boxed{
u(t)
=
u_k,
\qquad
t_k\leq t<t_{k+1}.
}
\]

The decision variables become

\[
\boxed{
u_0,u_1,\ldots,u_{N-1}.
}
\]

The state equations are integrated numerically.

%%%%%%%%%%%%%%%%%%%%%%%%%%%%%%%%%%%%%%%%%%%%%%%%%%%%%%%%%%%%
\subsection{Direct Transcription}
\label{subsec:direct-transcription}
%%%%%%%%%%%%%%%%%%%%%%%%%%%%%%%%%%%%%%%%%%%%%%%%%%%%%%%%%%%%

Direct transcription discretizes both states and controls.

For example, Euler discretization gives

\[
\boxed{
x_{k+1}
=
x_k
+
h
f(x_k,u_k,t_k).
}
\]

These equations become equality constraints in a finite-dimensional
optimization problem.

%%%%%%%%%%%%%%%%%%%%%%%%%%%%%%%%%%%%%%%%%%%%%%%%%%%%%%%%%%%%
\subsection{Trapezoidal Transcription}
\label{subsec:trapezoidal-transcription}
%%%%%%%%%%%%%%%%%%%%%%%%%%%%%%%%%%%%%%%%%%%%%%%%%%%%%%%%%%%%

A more accurate approximation is

\[
\boxed{
x_{k+1}
-
x_k
=
\frac{h}{2}
\left[
f(x_k,u_k,t_k)
+
f(x_{k+1},u_{k+1},t_{k+1})
\right].
}
\]

This is based on the trapezoidal integration rule.

%%%%%%%%%%%%%%%%%%%%%%%%%%%%%%%%%%%%%%%%%%%%%%%%%%%%%%%%%%%%
\subsection{Collocation}
\label{subsec:optimal-control-collocation}
%%%%%%%%%%%%%%%%%%%%%%%%%%%%%%%%%%%%%%%%%%%%%%%%%%%%%%%%%%%%

Collocation approximates the state trajectory with polynomial functions
and enforces the differential equation at selected collocation points.

High-order collocation methods can achieve accurate trajectories with
relatively few mesh intervals.

They are widely used in practical trajectory optimization.

%%%%%%%%%%%%%%%%%%%%%%%%%%%%%%%%%%%%%%%%%%%%%%%%%%%%%%%%%%%%
\subsection{Direct Versus Indirect Methods}
\label{subsec:direct-versus-indirect}
%%%%%%%%%%%%%%%%%%%%%%%%%%%%%%%%%%%%%%%%%%%%%%%%%%%%%%%%%%%%

Indirect methods:

\[
\boxed{
\text{derive optimality equations first}.
}
\]

Direct methods:

\[
\boxed{
\text{discretize first and solve numerically}.
}
\]

Indirect methods can provide high accuracy and analytical insight.

Direct methods are usually easier to apply to complicated constrained
problems.

%%%%%%%%%%%%%%%%%%%%%%%%%%%%%%%%%%%%%%%%%%%%%%%%%%%%%%%%%%%%
\subsection{Mesh Refinement}
\label{subsec:optimal-control-mesh-refinement}
%%%%%%%%%%%%%%%%%%%%%%%%%%%%%%%%%%%%%%%%%%%%%%%%%%%%%%%%%%%%

A direct solution depends on the discretization mesh.

After solving, one may:

\begin{enumerate}
    \item estimate discretization error;
    \item refine intervals with large error;
    \item resolve the optimization problem;
    \item repeat until accuracy is acceptable.
\end{enumerate}

This is called mesh refinement.

%%%%%%%%%%%%%%%%%%%%%%%%%%%%%%%%%%%%%%%%%%%%%%%%%%%%%%%%%%%%
\subsection{Control Saturation}
\label{subsec:control-saturation}
%%%%%%%%%%%%%%%%%%%%%%%%%%%%%%%%%%%%%%%%%%%%%%%%%%%%%%%%%%%%

A computed unconstrained control may exceed physical actuator limits.

For example,

\[
u^*_{\mathrm{free}}
=
-Kx
\]

may violate

\[
u_{\min}
\leq
u
\leq
u_{\max}.
\]

Simply clipping the control may not preserve global optimality.

The bounds should ideally be included in the optimization problem itself.

%%%%%%%%%%%%%%%%%%%%%%%%%%%%%%%%%%%%%%%%%%%%%%%%%%%%%%%%%%%%
\subsection{State Estimation and Control}
\label{subsec:state-estimation-control}
%%%%%%%%%%%%%%%%%%%%%%%%%%%%%%%%%%%%%%%%%%%%%%%%%%%%%%%%%%%%

Feedback control requires knowledge of the state.

If the entire state is not measured, an estimator may be used.

This leads to the combination of:

\[
\boxed{
\text{state estimation}
+
\text{optimal control}.
}
\]

For linear Gaussian systems, the Kalman filter and LQR combine into the
Linear-Quadratic-Gaussian framework under standard assumptions.

%%%%%%%%%%%%%%%%%%%%%%%%%%%%%%%%%%%%%%%%%%%%%%%%%%%%%%%%%%%%
\subsection{Stochastic Optimal Control}
\label{subsec:stochastic-optimal-control}
%%%%%%%%%%%%%%%%%%%%%%%%%%%%%%%%%%%%%%%%%%%%%%%%%%%%%%%%%%%%

Suppose dynamics include random disturbances:

\[
\boxed{
dX_t
=
b
(
X_t,
u_t,t
)
\,dt
+
\sigma
(
X_t,u_t,t
)
\,dW_t,
}
\]

where \(W_t\) is Brownian motion.

The objective may be

\[
\boxed{
\min
\mathbb{E}
\left[
\Phi(X_T)
+
\int_t^T
L(X_s,u_s,s)\,ds
\right].
}
\]

This is a stochastic optimal-control problem.

%%%%%%%%%%%%%%%%%%%%%%%%%%%%%%%%%%%%%%%%%%%%%%%%%%%%%%%%%%%%
\subsection{Stochastic HJB Equation}
\label{subsec:stochastic-hjb}
%%%%%%%%%%%%%%%%%%%%%%%%%%%%%%%%%%%%%%%%%%%%%%%%%%%%%%%%%%%%

Under appropriate smoothness assumptions, the value function can satisfy

\[
\boxed{
-
V_t
=
\min_u
\left\{
L
+
\nabla V^T b
+
\frac12
\operatorname{tr}
\left(
\sigma\sigma^T
\nabla^2V
\right)
\right\}.
}
\]

The second-derivative term arises from the diffusion process.

%%%%%%%%%%%%%%%%%%%%%%%%%%%%%%%%%%%%%%%%%%%%%%%%%%%%%%%%%%%%
\subsection{Deterministic HJB as a Special Case}
\label{subsec:deterministic-hjb-special-case}
%%%%%%%%%%%%%%%%%%%%%%%%%%%%%%%%%%%%%%%%%%%%%%%%%%%%%%%%%%%%

If

\[
\boxed{
\sigma=0,
}
\]

the diffusion term disappears.

Then the stochastic HJB equation reduces to

\[
\boxed{
-
V_t
=
\min_u
\left\{
L
+
\nabla V^Tf
\right\},
}
\]

which is the deterministic HJB equation.

%%%%%%%%%%%%%%%%%%%%%%%%%%%%%%%%%%%%%%%%%%%%%%%%%%%%%%%%%%%%
\subsection{Linear-Quadratic-Gaussian Control Preview}
\label{subsec:lqg-preview}
%%%%%%%%%%%%%%%%%%%%%%%%%%%%%%%%%%%%%%%%%%%%%%%%%%%%%%%%%%%%

Linear-Quadratic-Gaussian control combines:

\[
\boxed{
\text{linear dynamics},
}
\]

\[
\boxed{
\text{quadratic costs},
}
\]

and

\[
\boxed{
\text{Gaussian disturbances and measurement noise}.
}
\]

Under standard assumptions, optimal state estimation and optimal linear
feedback can be designed separately.

This is known as the separation principle.

%%%%%%%%%%%%%%%%%%%%%%%%%%%%%%%%%%%%%%%%%%%%%%%%%%%%%%%%%%%%
\subsection{Optimal Control and Reinforcement Learning}
\label{subsec:optimal-control-reinforcement-learning}
%%%%%%%%%%%%%%%%%%%%%%%%%%%%%%%%%%%%%%%%%%%%%%%%%%%%%%%%%%%%

Dynamic programming in optimal control and reinforcement learning share
the Bellman principle.

Optimal control usually begins with a known mathematical model.

Reinforcement learning may instead learn value functions or policies from
interaction or data.

The mathematical bridge includes:

\[
\boxed{
\text{value functions},
}
\]

\[
\boxed{
\text{Bellman equations},
}
\]

and

\[
\boxed{
\text{policies}.
}
\]

%%%%%%%%%%%%%%%%%%%%%%%%%%%%%%%%%%%%%%%%%%%%%%%%%%%%%%%%%%%%
\subsection{Policy Iteration Preview}
\label{subsec:policy-iteration-control}
%%%%%%%%%%%%%%%%%%%%%%%%%%%%%%%%%%%%%%%%%%%%%%%%%%%%%%%%%%%%

Policy iteration alternates between:

\[
\boxed{
\text{policy evaluation}
}
\]

and

\[
\boxed{
\text{policy improvement}.
}
\]

Given policy \(\pi\), first evaluate its value.

Then choose controls that improve the Bellman objective.

Repeat until the policy stabilizes.

%%%%%%%%%%%%%%%%%%%%%%%%%%%%%%%%%%%%%%%%%%%%%%%%%%%%%%%%%%%%
\subsection{Value Iteration Preview}
\label{subsec:value-iteration-control}
%%%%%%%%%%%%%%%%%%%%%%%%%%%%%%%%%%%%%%%%%%%%%%%%%%%%%%%%%%%%

Value iteration repeatedly applies a Bellman operator.

For discrete time,

\[
\boxed{
V_{k+1}(x)
=
\min_u
\left[
L(x,u)
+
\gamma
V_k
(
f(x,u)
)
\right]
}
\]

in a discounted deterministic setting.

The method seeks the fixed point of the Bellman equation.

%%%%%%%%%%%%%%%%%%%%%%%%%%%%%%%%%%%%%%%%%%%%%%%%%%%%%%%%%%%%
\subsection{Optimal Control and Optimization}
\label{subsec:optimal-control-optimization-connection}
%%%%%%%%%%%%%%%%%%%%%%%%%%%%%%%%%%%%%%%%%%%%%%%%%%%%%%%%%%%%

Ordinary optimization selects a finite-dimensional variable:

\[
\boxed{
x\in\mathbb{R}^n.
}
\]

Optimal control selects an entire function:

\[
\boxed{
u(\cdot).
}
\]

Thus optimal control is an infinite-dimensional optimization problem with
differential constraints.

%%%%%%%%%%%%%%%%%%%%%%%%%%%%%%%%%%%%%%%%%%%%%%%%%%%%%%%%%%%%
\subsection{Optimal Control and Dynamic Programming}
\label{subsec:optimal-control-dp-connection}
%%%%%%%%%%%%%%%%%%%%%%%%%%%%%%%%%%%%%%%%%%%%%%%%%%%%%%%%%%%%

Dynamic programming converts a global trajectory problem into local
recursive decisions.

The state summarizes all information relevant to future optimization.

The value function summarizes the optimal future cost.

Thus:

\[
\boxed{
\text{state}
+
\text{Bellman principle}
+
\text{value function}
}
\]

form the central dynamic-programming view of optimal control.

%%%%%%%%%%%%%%%%%%%%%%%%%%%%%%%%%%%%%%%%%%%%%%%%%%%%%%%%%%%%
\subsection{Optimal Control and Nonlinear Programming}
\label{subsec:optimal-control-nlp-connection}
%%%%%%%%%%%%%%%%%%%%%%%%%%%%%%%%%%%%%%%%%%%%%%%%%%%%%%%%%%%%

After direct transcription,

\[
\boxed{
\text{states}
+
\text{controls}
}
\]

at discretization points become ordinary optimization variables.

Dynamics become equality constraints.

Path constraints become inequalities.

Thus nonlinear-programming algorithms can solve large optimal-control
problems numerically.

%%%%%%%%%%%%%%%%%%%%%%%%%%%%%%%%%%%%%%%%%%%%%%%%%%%%%%%%%%%%
\subsection{Optimal Control and Differential Equations}
\label{subsec:optimal-control-ode-connection}
%%%%%%%%%%%%%%%%%%%%%%%%%%%%%%%%%%%%%%%%%%%%%%%%%%%%%%%%%%%%

The state dynamics determine what trajectories are physically possible.

Optimal-control algorithms repeatedly solve or approximate differential
equations.

Therefore numerical ODE accuracy can directly affect optimization
accuracy.

This creates a strong connection with Chapter~9 and Chapter~11.

%%%%%%%%%%%%%%%%%%%%%%%%%%%%%%%%%%%%%%%%%%%%%%%%%%%%%%%%%%%%
\subsection{Optimal Control Workflow}
\label{subsec:optimal-control-workflow}
%%%%%%%%%%%%%%%%%%%%%%%%%%%%%%%%%%%%%%%%%%%%%%%%%%%%%%%%%%%%

A practical workflow is:

\begin{enumerate}
    \item identify the state variables;
    \item identify the control variables;
    \item construct the dynamical equations;
    \item specify initial conditions;
    \item specify terminal conditions;
    \item define running and terminal costs;
    \item specify control bounds;
    \item specify state and path constraints;
    \item determine whether terminal time is fixed or free;
    \item check reachability or controllability when appropriate;
    \item determine whether analytical optimality conditions are useful;
    \item construct the Hamiltonian for an indirect approach;
    \item derive state, costate, stationarity, and transversality
          conditions;
    \item alternatively formulate the Bellman or HJB equation;
    \item choose direct or indirect numerical methods;
    \item solve the resulting boundary-value or nonlinear optimization
          problem;
    \item verify dynamical feasibility;
    \item inspect active constraints and switching behavior;
    \item test sensitivity to model parameters;
    \item simulate the resulting control under perturbed conditions;
    \item implement feedback when robustness to disturbances is needed.
\end{enumerate}

%%%%%%%%%%%%%%%%%%%%%%%%%%%%%%%%%%%%%%%%%%%%%%%%%%%%%%%%%%%%
\subsection{Choosing an Optimal-Control Method}
\label{subsec:choosing-optimal-control-method}
%%%%%%%%%%%%%%%%%%%%%%%%%%%%%%%%%%%%%%%%%%%%%%%%%%%%%%%%%%%%

For analytical necessary conditions, consider:

\[
\boxed{
\text{Pontryagin's Maximum Principle}.
}
\]

For small-state feedback problems, consider:

\[
\boxed{
\text{dynamic programming and HJB}.
}
\]

For linear systems with quadratic objectives, use:

\[
\boxed{
\text{LQR}.
}
\]

For constrained real-time control, consider:

\[
\boxed{
\text{MPC}.
}
\]

For complicated nonlinear constrained trajectories, consider:

\[
\boxed{
\text{direct transcription or collocation}.
}
\]

For smooth boundary-value problems with good initial guesses, consider:

\[
\boxed{
\text{shooting methods}.
}
\]

%%%%%%%%%%%%%%%%%%%%%%%%%%%%%%%%%%%%%%%%%%%%%%%%%%%%%%%%%%%%
\subsection{Common Errors}
\label{subsec:optimal-control-common-errors}
%%%%%%%%%%%%%%%%%%%%%%%%%%%%%%%%%%%%%%%%%%%%%%%%%%%%%%%%%%%%

\subsubsection{Treating State and Control as Independent Variables}

The state must satisfy

\[
\boxed{
\dot{x}=f(x,u,t).
}
\]

One cannot optimize \(x(t)\) independently of the dynamics.

%%%%%%%%%%%%%%%%%%%%%%%%%%%%%%%%%%%%%%%%%%%%%%%%%%%%%%%%%%%%
\subsubsection{Using Inconsistent Hamiltonian Signs}

If using

\[
H=L+\lambda^Tf
\]

for minimization, minimize the Hamiltonian with respect to the control.

If using another convention, all costate and extremum conditions must be
changed consistently.

%%%%%%%%%%%%%%%%%%%%%%%%%%%%%%%%%%%%%%%%%%%%%%%%%%%%%%%%%%%%
\subsubsection{Forgetting the Costate Equation}

The state equation alone does not produce Pontryagin optimality.

The coupled system includes

\[
\boxed{
\dot{\lambda}
=
-H_x.
}
\]

%%%%%%%%%%%%%%%%%%%%%%%%%%%%%%%%%%%%%%%%%%%%%%%%%%%%%%%%%%%%
\subsubsection{Forgetting Terminal Conditions}

The state usually has an initial condition.

The costate often has a terminal condition.

Without the correct endpoint conditions, the boundary-value problem is
incomplete.

%%%%%%%%%%%%%%%%%%%%%%%%%%%%%%%%%%%%%%%%%%%%%%%%%%%%%%%%%%%%
\subsubsection{Using \(H_u=0\) With a Boundary Control}

The condition

\[
H_u=0
\]

applies to suitable interior controls.

With bounds, the optimum may lie at

\[
u_{\min}
\]

or

\[
u_{\max}.
\]

%%%%%%%%%%%%%%%%%%%%%%%%%%%%%%%%%%%%%%%%%%%%%%%%%%%%%%%%%%%%
\subsubsection{Assuming Bang-Bang Control for Every Bounded Problem}

Bang-bang behavior commonly arises when the Hamiltonian is affine in the
control.

Quadratic control penalties typically produce interior controls instead.

%%%%%%%%%%%%%%%%%%%%%%%%%%%%%%%%%%%%%%%%%%%%%%%%%%%%%%%%%%%%
\subsubsection{Ignoring Singular Controls}

If the switching function vanishes over an interval, ordinary bang-bang
logic may fail.

Additional analysis is required.

%%%%%%%%%%%%%%%%%%%%%%%%%%%%%%%%%%%%%%%%%%%%%%%%%%%%%%%%%%%%
\subsubsection{Confusing PMP Conditions With a Global Optimality Proof}

Pontryagin conditions are generally necessary conditions.

A trajectory satisfying them need not automatically be the global optimum.

Convexity or additional verification conditions may be needed.

%%%%%%%%%%%%%%%%%%%%%%%%%%%%%%%%%%%%%%%%%%%%%%%%%%%%%%%%%%%%
\subsubsection{Confusing Open-Loop and Feedback Solutions}

A precomputed function \(u(t)\) does not react directly to state errors.

A feedback law

\[
u=\pi(x,t)
\]

does.

%%%%%%%%%%%%%%%%%%%%%%%%%%%%%%%%%%%%%%%%%%%%%%%%%%%%%%%%%%%%
\subsubsection{Ignoring Controllability}

If the target state cannot be reached under the dynamics, no optimization
algorithm can produce a feasible trajectory satisfying the target exactly.

%%%%%%%%%%%%%%%%%%%%%%%%%%%%%%%%%%%%%%%%%%%%%%%%%%%%%%%%%%%%
\subsubsection{Ignoring Units in Quadratic Costs}

Terms such as

\[
x^TQx
\]

and

\[
u^TRu
\]

may use different physical units.

The weights \(Q\) and \(R\) should be chosen with scaling and physical
interpretation in mind.

%%%%%%%%%%%%%%%%%%%%%%%%%%%%%%%%%%%%%%%%%%%%%%%%%%%%%%%%%%%%
\subsubsection{Assuming Larger LQR Weight Always Improves Performance}

Increasing a state penalty can require much larger control effort.

Actuator constraints and model error may make aggressive gains
undesirable.

%%%%%%%%%%%%%%%%%%%%%%%%%%%%%%%%%%%%%%%%%%%%%%%%%%%%%%%%%%%%
\subsubsection{Ignoring Control Saturation}

An unconstrained LQR solution may request physically impossible control.

Actuator limits should be included when they materially affect operation.

%%%%%%%%%%%%%%%%%%%%%%%%%%%%%%%%%%%%%%%%%%%%%%%%%%%%%%%%%%%%
\subsubsection{Confusing HJB With an Ordinary Differential Equation}

For multidimensional continuous-state problems, HJB is generally a
nonlinear partial differential equation.

Its numerical solution can be difficult.

%%%%%%%%%%%%%%%%%%%%%%%%%%%%%%%%%%%%%%%%%%%%%%%%%%%%%%%%%%%%
\subsubsection{Ignoring the Curse of Dimensionality}

A grid-based dynamic-programming method may become impractical as state
dimension grows.

Alternative approximations may be required.

%%%%%%%%%%%%%%%%%%%%%%%%%%%%%%%%%%%%%%%%%%%%%%%%%%%%%%%%%%%%
\subsubsection{Using Too Coarse a Direct-Transcription Mesh}

A discretized trajectory can satisfy the numerical NLP while poorly
approximating the true continuous dynamics.

Mesh refinement and simulation verification are important.

%%%%%%%%%%%%%%%%%%%%%%%%%%%%%%%%%%%%%%%%%%%%%%%%%%%%%%%%%%%%
\subsubsection{Assuming a Numerical Solver's Success Flag Proves Physical Feasibility}

The optimized trajectory should be independently checked by integrating
the original dynamics and evaluating all constraints.

%%%%%%%%%%%%%%%%%%%%%%%%%%%%%%%%%%%%%%%%%%%%%%%%%%%%%%%%%%%%
\subsubsection{Ignoring Model Error}

An open-loop optimum for an inaccurate model may perform poorly in the
physical system.

Feedback, robust control, stochastic control, or repeated reoptimization
may be needed.

%%%%%%%%%%%%%%%%%%%%%%%%%%%%%%%%%%%%%%%%%%%%%%%%%%%%%%%%%%%%
\subsection{Applications}
\label{subsec:optimal-control-applications}
%%%%%%%%%%%%%%%%%%%%%%%%%%%%%%%%%%%%%%%%%%%%%%%%%%%%%%%%%%%%

\begin{applicationbox}

\textbf{Aerospace.}
Optimal control determines launch, spacecraft, aircraft, descent, and
minimum-fuel trajectories.

\medskip

\textbf{Robotics.}
Robots optimize motion while respecting nonlinear dynamics, actuator
limits, contact constraints, and obstacle avoidance.

\medskip

\textbf{Automotive Systems.}
Optimal control is used for trajectory tracking, autonomous driving,
energy management, and vehicle stability.

\medskip

\textbf{Energy Systems.}
Batteries, generators, storage systems, and grids require time-dependent
control under physical and economic constraints.

\medskip

\textbf{Economics.}
Consumption, investment, resource extraction, and capital accumulation can
be formulated as intertemporal control problems.

\medskip

\textbf{Biology and Medicine.}
Intervention schedules can be optimized for epidemics, drug dosing,
population management, and treatment planning.

\medskip

\textbf{Manufacturing.}
Production rates, temperatures, pressures, and process conditions can be
optimized dynamically.

\medskip

\textbf{Finance.}
Portfolio and consumption decisions over continuous time lead to
stochastic optimal-control models.

\end{applicationbox}

%%%%%%%%%%%%%%%%%%%%%%%%%%%%%%%%%%%%%%%%%%%%%%%%%%%%%%%%%%%%
\subsection{Exercises}
\label{subsec:optimal-control-exercises}
%%%%%%%%%%%%%%%%%%%%%%%%%%%%%%%%%%%%%%%%%%%%%%%%%%%%%%%%%%%%

\begin{exercise}[1]
Define a state variable.
\end{exercise}

\begin{exercise}[1]
Define a control variable.
\end{exercise}

\begin{exercise}[1]
Define running cost.
\end{exercise}

\begin{exercise}[1]
Define terminal cost.
\end{exercise}

\begin{exercise}[1]
Define an admissible control.
\end{exercise}

\begin{exercise}[1]
Distinguish open-loop and feedback control.
\end{exercise}

\begin{exercise}[1]
Write the general Bolza optimal-control problem.
\end{exercise}

\begin{exercise}[1]
Define the Hamiltonian using the minimization convention of this section.
\end{exercise}

\begin{exercise}[1]
State the costate equation.
\end{exercise}

\begin{exercise}[1]
State the HJB equation.
\end{exercise}

\begin{exercise}[2]
For

\[
\dot{x}=ax+bu,
\]

and running cost

\[
L=\frac12(qx^2+ru^2),
\]

construct the Hamiltonian.
\end{exercise}

\begin{exercise}[2]
Differentiate the previous Hamiltonian with respect to \(u\).
\end{exercise}

\begin{exercise}[2]
Solve the resulting stationarity equation for \(u\).
\end{exercise}

\begin{exercise}[2]
Write the costate equation for the same problem.
\end{exercise}

\begin{exercise}[2]
For terminal cost

\[
\Phi(x(T))
=
\frac12Fx(T)^2,
\]

determine the costate terminal condition.
\end{exercise}

\begin{exercise}[3]
Solve the problem

\[
\dot{x}=u,
\qquad
x(0)=x_0,
\]

with

\[
J
=
\frac12x(T)^2
+
\frac12
\int_0^Tu^2dt
\]

using Pontryagin's conditions.
\end{exercise}

\begin{exercise}[3]
Verify the optimal state and control obtained in the previous problem.
\end{exercise}

\begin{exercise}[3]
Construct the Hamiltonian for a minimum-time problem.
\end{exercise}

\begin{exercise}[3]
Explain why bounded minimum-time controls often have bang-bang structure.
\end{exercise}

\begin{exercise}[3]
Given switching function

\[
\psi(t)=2-t
\]

and control bounds

\[
-1\leq u\leq1,
\]

determine the minimizing bang-bang control away from the switching time.
\end{exercise}

\begin{exercise}[3]
Identify the switching time in the previous problem.
\end{exercise}

\begin{exercise}[3]
Explain what a singular arc means.
\end{exercise}

\begin{exercise}[3]
Construct the Hamiltonian for the finite-horizon LQR problem.
\end{exercise}

\begin{exercise}[3]
Derive

\[
u^*
=
-R^{-1}B^T\lambda.
\]
\end{exercise}

\begin{exercise}[3]
Derive the LQR costate equation.
\end{exercise}

\begin{exercise}[3]
Assume

\[
\lambda=Px
\]

and derive the Riccati differential equation.
\end{exercise}

\begin{exercise}[3]
Write the finite-horizon LQR feedback law.
\end{exercise}

\begin{exercise}[3]
Explain the roles of \(Q\) and \(R\) in LQR.
\end{exercise}

\begin{exercise}[3]
Write the algebraic Riccati equation for infinite-horizon continuous-time
LQR.
\end{exercise}

\begin{exercise}[3]
Construct the controllability matrix for a two-dimensional linear system.
\end{exercise}

\begin{exercise}[3]
Determine whether the system is controllable by calculating its rank.
\end{exercise}

\begin{exercise}[3]
Write a discrete-time finite-horizon optimal-control problem.
\end{exercise}

\begin{exercise}[3]
Derive its Bellman recursion.
\end{exercise}

\begin{exercise}[3]
Explain how Bellman recursion produces feedback.
\end{exercise}

\begin{exercise}[3]
Write the deterministic continuous-time HJB equation.
\end{exercise}

\begin{exercise}[3]
For

\[
\dot{x}=u,
\qquad
L=\frac12u^2,
\]

find the HJB minimizing control in terms of \(V_x\).
\end{exercise}

\begin{exercise}[3]
Explain the relationship between the costate and value-function gradient.
\end{exercise}

\begin{exercise}[3]
Compare PMP and HJB approaches.
\end{exercise}

\begin{exercise}[3]
Construct a bounded state constraint.
\end{exercise}

\begin{exercise}[3]
Construct a bounded control constraint.
\end{exercise}

\begin{exercise}[3]
Write an augmented Hamiltonian for a path inequality constraint.
\end{exercise}

\begin{exercise}[3]
Write the corresponding complementarity conditions.
\end{exercise}

\begin{exercise}[3]
Construct a quadratic trajectory-tracking objective.
\end{exercise}

\begin{exercise}[3]
Construct an MPC objective for a discrete-time linear system.
\end{exercise}

\begin{exercise}[3]
Explain the receding-horizon principle.
\end{exercise}

\begin{exercise}[3]
Explain why MPC produces feedback despite solving an open-loop trajectory
at each update.
\end{exercise}

\begin{exercise}[3]
Discretize

\[
\dot{x}=f(x,u,t)
\]

using forward Euler.
\end{exercise}

\begin{exercise}[3]
Discretize the same dynamics using the trapezoidal rule.
\end{exercise}

\begin{exercise}[3]
Explain direct transcription.
\end{exercise}

\begin{exercise}[3]
Explain the shooting method.
\end{exercise}

\begin{exercise}[3]
Compare direct and indirect numerical methods.
\end{exercise}

\begin{exercise}[3]
Explain why mesh refinement may be necessary.
\end{exercise}

\begin{exercise}[3]
Write a stochastic controlled diffusion.
\end{exercise}

\begin{exercise}[3]
Identify the additional second-order term in the stochastic HJB equation.
\end{exercise}

\begin{exercise}[3]
Explain the relationship between optimal control and reinforcement
learning.
\end{exercise}

\begin{exercise}[4]
Derive Pontryagin's necessary conditions using constrained variational
arguments.
\end{exercise}

\begin{exercise}[4]
Study transversality conditions for free terminal states.
\end{exercise}

\begin{exercise}[4]
Study transversality conditions for free terminal time.
\end{exercise}

\begin{exercise}[4]
Study abnormal extremals in Pontryagin's Maximum Principle.
\end{exercise}

\begin{exercise}[4]
Study sufficient optimality conditions for convex optimal-control
problems.
\end{exercise}

\begin{exercise}[4]
Investigate bang-bang optimal control for a double integrator.
\end{exercise}

\begin{exercise}[4]
Derive the switching curve for minimum-time double-integrator control.
\end{exercise}

\begin{exercise}[4]
Study singular optimal controls.
\end{exercise}

\begin{exercise}[4]
Study the Legendre--Clebsch condition.
\end{exercise}

\begin{exercise}[4]
Derive finite-horizon LQR using dynamic programming instead of PMP.
\end{exercise}

\begin{exercise}[4]
Show that the HJB quadratic ansatz produces the Riccati equation.
\end{exercise}

\begin{exercise}[4]
Study stabilizability and detectability conditions for infinite-horizon
LQR.
\end{exercise}

\begin{exercise}[4]
Study discrete-time LQR and the discrete Riccati equation.
\end{exercise}

\begin{exercise}[4]
Study the Hamilton--Jacobi--Bellman verification theorem.
\end{exercise}

\begin{exercise}[4]
Study viscosity solutions of HJB equations conceptually.
\end{exercise}

\begin{exercise}[4]
Investigate the curse of dimensionality in numerical dynamic programming.
\end{exercise}

\begin{exercise}[4]
Study state-constrained Pontryagin conditions.
\end{exercise}

\begin{exercise}[4]
Study direct collocation methods.
\end{exercise}

\begin{exercise}[4]
Study pseudospectral optimal-control methods.
\end{exercise}

\begin{exercise}[4]
Investigate multiple shooting for nonlinear trajectory optimization.
\end{exercise}

\begin{exercise}[4]
Study sequential quadratic programming for direct transcription.
\end{exercise}

\begin{exercise}[4]
Study constrained linear MPC.
\end{exercise}

\begin{exercise}[4]
Investigate terminal sets and stability in MPC.
\end{exercise}

\begin{exercise}[4]
Study stochastic HJB equations.
\end{exercise}

\begin{exercise}[4]
Study Linear-Quadratic-Gaussian control and the separation principle.
\end{exercise}

\begin{exercise}[4]
Investigate risk-sensitive optimal control.
\end{exercise}

\begin{exercise}[4]
Study robust optimal control.
\end{exercise}

\begin{exercise}[4]
Investigate differential games as a connection between optimal control and
game theory.
\end{exercise}

%%%%%%%%%%%%%%%%%%%%%%%%%%%%%%%%%%%%%%%%%%%%%%%%%%%%%%%%%%%%
\subsection{Conceptual Summary}
\label{subsec:optimal-control-summary}
%%%%%%%%%%%%%%%%%%%%%%%%%%%%%%%%%%%%%%%%%%%%%%%%%%%%%%%%%%%%

Optimal control optimizes decisions distributed over time while respecting
dynamical equations.

A general problem is

\[
\boxed{
\min_{u(\cdot)}
\left[
\Phi(x(T),T)
+
\int_{t_0}^{T}
L(x,u,t)\,dt
\right]
}
\]

subject to

\[
\boxed{
\dot{x}=f(x,u,t).
}
\]

Using the minimization convention,

\[
\boxed{
H
=
L+\lambda^Tf.
}
\]

Pontryagin's necessary conditions include

\[
\boxed{
\dot{x}
=
H_\lambda,
}
\]

\[
\boxed{
\dot{\lambda}
=
-H_x,
}
\]

and

\[
\boxed{
u^*
\in
\operatorname*{arg\,min}_{u\in U}H.
}
\]

For an unconstrained interior control,

\[
\boxed{
H_u=0.
}
\]

Dynamic programming introduces the value function and HJB equation:

\[
\boxed{
-V_t
=
\min_u
\left\{
L
+
\nabla V^Tf
\right\}.
}
\]

For linear-quadratic control,

\[
\boxed{
\dot{x}=Ax+Bu
}
\]

and

\[
\boxed{
J
=
\frac12x(T)^TFx(T)
+
\frac12
\int
(
x^TQx+u^TRu
)dt,
}
\]

the optimal feedback is

\[
\boxed{
u^*
=
-
R^{-1}B^TPx,
}
\]

where \(P\) satisfies a Riccati equation.

Modern optimal control therefore combines

\[
\boxed{
\text{optimization}
+
\text{differential equations}
+
\text{dynamic programming}
+
\text{numerical computation}
+
\text{feedback}.
}
\]

%%%%%%%%%%%%%%%%%%%%%%%%%%%%%%%%%%%%%%%%%%%%%%%%%%%%%%%%%%%%
\subsection{Section Connections}
\label{subsec:optimal-control-connections}
%%%%%%%%%%%%%%%%%%%%%%%%%%%%%%%%%%%%%%%%%%%%%%%%%%%%%%%%%%%%

\chapterconnection{
Optimal Control completes Chapter~10 by extending optimization from
finite-dimensional decision vectors to time-dependent decisions governed
by dynamical systems. Pontryagin's Maximum Principle generalizes
Lagrange-multiplier reasoning to differential constraints, while the
Hamilton--Jacobi--Bellman equation connects Optimal Control directly with
Dynamic Programming. Linear-quadratic regulation provides one of the most
important analytically solvable control families, and direct transcription
connects modern trajectory optimization with nonlinear programming.
Stochastic Optimal Control extends these ideas using the stochastic
processes of Chapter~7. The next chapter, Numerical and Computational
Mathematics, develops the computational methods required to solve many of
the nonlinear equations, differential equations, optimization problems,
Monte Carlo models, and large numerical systems encountered throughout
advanced mathematics.
}

%%%%%%%%%%%%%%%%%%%%%%%%%%%%%%%%%%%%%%%%%%%%%%%%%%%%%%%%%%%%
% SECTION SOURCE TRACKING
%%%%%%%%%%%%%%%%%%%%%%%%%%%%%%%%%%%%%%%%%%%%%%%%%%%%%%%%%%%%

%------------------------------------------------------------
% 10.16 Optimal Control
%------------------------------------------------------------
%
% Source basis:
%   - Standard deterministic optimal-control theory.
%   - Standard dynamic-programming and Hamilton--Jacobi--Bellman theory.
%   - Standard linear-quadratic control.
%   - Standard numerical optimal-control methods.
%
% Used for:
%   - states and controls;
%   - dynamical constraints;
%   - admissible controls;
%   - Bolza, Lagrange, and Mayer forms;
%   - running and terminal costs;
%   - fixed and free terminal conditions;
%   - open-loop and feedback control;
%   - calculus-of-variations connections;
%   - Hamiltonian formulation;
%   - costates;
%   - Pontryagin's Maximum Principle;
%   - sign conventions;
%   - state and costate equations;
%   - transversality;
%   - pointwise Hamiltonian minimization;
%   - control constraints;
%   - two-point boundary-value problems;
%   - free-final-time conditions;
%   - minimum-time control;
%   - bang-bang control;
%   - switching functions;
%   - singular arcs;
%   - minimum-energy control;
%   - trajectory tracking;
%   - LQR;
%   - Riccati differential equations;
%   - infinite-horizon LQR;
%   - controllability;
%   - discrete-time optimal control;
%   - Bellman recursion;
%   - value functions;
%   - HJB equations;
%   - PMP versus HJB;
%   - curse of dimensionality;
%   - state, control, path, and endpoint constraints;
%   - terminal penalties;
%   - discounted control;
%   - economic and biological control examples;
%   - trajectory optimization;
%   - MPC;
%   - direct and indirect numerical methods;
%   - shooting and multiple shooting;
%   - direct transcription;
%   - collocation;
%   - mesh refinement;
%   - control saturation;
%   - state estimation connections;
%   - stochastic optimal control;
%   - stochastic HJB;
%   - LQG preview;
%   - reinforcement-learning connections;
%   - applications and exercises.
%
% Notes:
%   - This completes Chapter 10 — Optimization and Operations Research.
%   - Chapter 11 begins with Numerical Analysis.
%   - Differential-equation foundations appear in Chapter 9.
%   - Stochastic-process foundations appear in Chapter 7.
%   - Numerical methods for ODEs, optimization, Monte Carlo simulation,
%     and scientific computing are developed in Chapter 11.
%
%%%%%%%%%%%%%%%%%%%%%%%%%%%%%%%%%%%%%%%%%%%%%%%%%%%%%%%%%%%%